% !Mode:: "TeX:UTF-8"

%%%%%%%%%% Package %%%%%%%%%%%%
% 项目已整体迁移到 XeLaTeX，中文由 ctex + xeCJK 处理
\usepackage{enumerate}
\usepackage[chapter]{algorithm}
\usepackage{algorithmic}
\usepackage{graphicx}                       % xelatex 下自动使用 xetex
% 驱动，直接支持 PDF/PNG/JPG
\usepackage{geometry}
\geometry{left=3.0cm,right=3.0cm,top=3.0cm,bottom=3.0cm}%%%%%%%页面边距的设置
% 支持版面尺寸设置
\usepackage{titlesec}                       % 控制标题的宏包
\usepackage{titletoc}                       % 控制目录的宏包
\usepackage{fancyhdr}                       % fancyhdr 宏包 支持页眉和页脚的相关定义
\usepackage{ctex}                          % xelatex 下自动加载 xeCJK，支持中文
% 将模板自定义的 \song/\hei/\kai/\fs/\li/\FSGB（\CJKfamily{...}）映射到 Windows 系统字体
\setCJKfamilyfont{song}{SimSun}
\setCJKfamilyfont{hei}{SimHei}
\setCJKfamilyfont{kai}{KaiTi}
\setCJKfamilyfont{fs}{FangSong}
\setCJKfamilyfont{li}{LiSu}
\setCJKfamilyfont{FSGB}{FangSong}
% 拉丁正文用 Times New Roman（fontspec）；数学用 newtxmath（Times 风格）。
% 必须显式先加载 fontspec，再 \setmainfont —— 否则 newtxmath 解析 Times 粗体(OT1/b/n)时
% 会在日志报 "Font shape OT1/TimesNewRoman(0)/b/n undefined"（虽不影响粗体渲染，但不干净）。
% 注意：不要在这里加载 newtxtext —— 它与 fontspec 冲突，会导致 SimSun 等中文字体解析失败。
\usepackage{fontspec}
\setmainfont{Times New Roman}
\usepackage{cite}
\usepackage{color}                          % 支持彩色
\usepackage{amsmath}                        % AMSLaTeX 宏包 用来排出更加漂亮的公式
\usepackage{amssymb}                        % 数学符号生成命令
\usepackage[below]{placeins}                %允许上一个 section
% 的浮动图形出现在下一个 section 的开始部分，还提供\FloatBarrier 命令，使所有未处理的浮动图形立即被处理
\usepackage{flafter}                        %
% 使得所有浮动体不能被放置在其浮动环境之前，以免浮动体在引述它的文本之前出现。
\usepackage{multirow}                       % 使用 Multirow 宏包，使得表格可以合并多个 row 格
\usepackage{booktabs}                       %
% 表格，横的粗线；\specialrule{1pt}{0pt}{0pt}
\usepackage{longtable}                      % 支持跨页的表格。
\usepackage{tabularx}                       % 自动设置表格的列宽
\usepackage{setspace}
\usepackage{subfigure}                      % 支持子图 %centerlast 设置最后一行是否居中
\usepackage[subfigure]{ccaption}            % 支持子图的中文标题
\usepackage[sort&compress,numbers]{natbib}  % 支持引用缩写的宏包
\usepackage{enumitem}                       % 使用 enumitem 宏包，改变列表项的格式
\usepackage{enumerate}
\usepackage{calc}                           % 长度可以用 + - * / 进行计算
\usepackage{newtxmath}                       % Times 风格数学字体（xelatex 兼容）
\usepackage{bm}                             % 处理数学公式中的黑斜体的宏包
\usepackage[amsmath,thmmarks,hyperref]{ntheorem}  % 定理类环境宏包，其中
% amsmath 选项用来兼容 AMS LaTeX 的宏包
\usepackage{indentfirst}                    % 首行缩进宏包
\usepackage{fancyhdr}
\usepackage{lastpage}
\usepackage{layout}
\usepackage[titles,subfigure]{tocloft}
%控制生成的表格和图片的目录格式
\usepackage{caption}

\usepackage{ragged2e}

%如果您的 pdf 制作中文书签有乱码使用如下命令，就可以解决了
% xelatex 下 hyperref 自动检测 xetex 驱动，无需显式指定
\usepackage[unicode,
  pdfstartview=FitH,
  bookmarksnumbered=true,
  bookmarksopen=true,
  colorlinks=true,
  pdfborder={0 0 1},
  citecolor=black,
  linkcolor=black,
  anchorcolor=black,
  urlcolor=black,
  breaklinks=true
]{hyperref}

\renewcommand{\theequation}{\arabic{chapter}-\arabic{equation}}
\makeatletter
\def\tagform@#1{\maketag@@@{（\ignorespaces#1\unskip）}}
\makeatother
